%!TEX program = xelatex
\documentclass[12pt,a4paper,oneside]{article}

% ============================================================================
% PACKAGES & CONFIGURATION
% ============================================================================
\usepackage[utf-8]{inputenc}
\usepackage[left=2.5cm,right=2.5cm,top=3cm,bottom=3cm]{geometry}
\usepackage{graphicx}
\usepackage{fancyhdr}
\usepackage{color}
\usepackage{listings}
\usepackage{hyperref}
\usepackage{xcolor}
\usepackage{tcolorbox}
\usepackage{array}
\usepackage{booktabs}
\usepackage{amsmath}
\usepackage{amssymb}
\usepackage{float}
\usepackage{titlesec}
\usepackage{enumitem}
\usepackage{fancyvrb}
\usepackage{longtable}

% Color definitions
\definecolor{codegreen}{rgb}{0.0, 0.5, 0.0}
\definecolor{codegray}{rgb}{0.5, 0.5, 0.5}
\definecolor{codepurple}{rgb}{0.58, 0, 0.82}
\definecolor{backcolour}{rgb}{0.95, 0.95, 0.95}
\definecolor{lightblue}{rgb}{0.9, 0.95, 1.0}
\definecolor{darkblue}{rgb}{0.0, 0.2, 0.6}
\definecolor{lightyellow}{rgb}{1.0, 0.98, 0.8}
\definecolor{lightgreen}{rgb}{0.9, 1.0, 0.9}
\definecolor{lightred}{rgb}{1.0, 0.9, 0.9}
\definecolor{orange}{rgb}{1.0, 0.65, 0.0}

% Code listing configuration
\lstset{
    backgroundcolor=\color{backcolour},
    commentstyle=\color{codegreen},
    keywordstyle=\color{codepurple},
    numberstyle=\tiny\color{codegray},
    stringstyle=\color{codegreen},
    basicstyle=\ttfamily\small,
    breakatwhitespace=false,
    breaklines=true,
    captionpos=b,
    keepspaces=true,
    numbers=left,
    numbersep=5pt,
    showspaces=false,
    showstringspaces=false,
    showtabs=false,
    tabsize=2,
    frame=single,
    language=json,
}

% Header and footer
\pagestyle{fancy}
\fancyhf{}
\rhead{\textcolor{darkblue}{\textit{VisioLearn Frontend Integration}}}
\lhead{\textcolor{darkblue}{API Reference}}
\cfoot{\thepage}
\renewcommand{\headrulewidth}{0.5pt}
\renewcommand{\footrulewidth}{0.5pt}

% Section styling
\titleformat{\section}{\Large\bfseries\color{darkblue}}{}{0pt}{}[\titlerule]
\titleformat{\subsection}{\large\bfseries\color{orange}}{}{0pt}{}
\titleformat{\subsubsection}{\bfseries\color{darkblue}}{}{0pt}{}

% ============================================================================
% TITLE PAGE
% ============================================================================
\title{%
    \vspace{2cm}
    \Huge\textbf{\textcolor{darkblue}{VisioLearn Frontend}} \\
    \vspace{0.5cm}
    \Large\textit{Backend API Integration Reference} \\
    \vspace{2cm}
    \large\textbf{Complete Backend API Documentation for Frontend Developers}
}

\author{%
    Backend API Team \\
    \vspace{0.5cm}
    Production API at \\
    \href{https://visiolearn-backend.onrender.com}{https://visiolearn-backend.onrender.com}
}

\date{\vspace{1cm}\large \textbf{Version 1.0} \\ May 2026}

% ============================================================================
% DOCUMENT
% ============================================================================
\begin{document}

% Title page
\maketitle
\thispagestyle{empty}
\newpage

% ============================================================================
% TABLE OF CONTENTS
% ============================================================================
\tableofcontents
\newpage

% ============================================================================
% SECTION 1: OVERVIEW & BASE CONFIGURATION
% ============================================================================
\section{Overview \& Configuration}

\subsection{API Base URL}

All requests must be sent to the following base URL:

\begin{tcolorbox}[colback=lightblue, colframe=darkblue, arc=0mm]
\texttt{https://visiolearn-backend.onrender.com/api/v1}
\end{tcolorbox}

\subsection{API Version}

This documentation covers \textbf{API v1}. All endpoints are versioned for backward compatibility.

\subsection{Supported HTTP Methods}

\begin{itemize}
    \item \textbf{GET}: Retrieve data
    \item \textbf{POST}: Create new data or perform actions
    \item \textbf{PUT}: Replace existing data
    \item \textbf{DELETE}: Remove data
\end{itemize}

\subsection{Content-Type}

All requests and responses use JSON format:
\begin{lstlisting}
Content-Type: application/json
\end{lstlisting}

For file uploads (voice, notes), use \texttt{multipart/form-data}.

\newpage

% ============================================================================
% SECTION 2: AUTHENTICATION
% ============================================================================
\section{Authentication \& Authorization}

\subsection{Overview}

VisioLearn uses JWT (JSON Web Token) authentication. All endpoints except login and signup require an authorization token.

\subsection{Authentication Flow}

\begin{enumerate}
    \item User provides email and password
    \item Backend returns \texttt{access\_token} and \texttt{refresh\_token}
    \item Frontend stores both tokens (preferably in secure storage)
    \item Frontend includes \texttt{access\_token} in every API request header
    \item When \texttt{access\_token} expires, use \texttt{refresh\_token} to get a new one
\end{enumerate}

\subsection{Token Storage Recommendations}

\begin{tcolorbox}[colback=lightgreen, colframe=darkblue]
\textbf{Security Note:} Store tokens securely:
\begin{itemize}
    \item \textbf{Best}: HttpOnly cookies (most secure, set by backend)
    \item \textbf{Good}: Secure local storage with encryption
    \item \textbf{Avoid}: Plain localStorage (vulnerable to XSS)
\end{itemize}
\end{tcolorbox}

\subsection{Authorization Header}

Include your access token in every authenticated request:

\begin{lstlisting}
Authorization: Bearer <access_token>
\end{lstlisting}

\subsection{User Roles}

The system supports three user roles with different permissions:

\begin{table}[H]
\centering
\small
\begin{tabular}{|l|l|}
\hline
\textbf{Role} & \textbf{Permissions} \\
\hline
\textbf{admin} & All endpoints, user management, school management \\
\hline
\textbf{teacher} & Upload notes, view student progress, manage lessons \\
\hline
\textbf{student} & Access lessons, initiate voice sessions, view progress \\
\hline
\end{tabular}
\caption{User Roles and Permissions}
\end{table}

\newpage

% ============================================================================
% SECTION 3: HTTP STATUS CODES
% ============================================================================
\section{HTTP Status Codes}

All responses include a status code. Understand what each means:

\begin{table}[H]
\centering
\small
\begin{tabular}{|c|l|}
\hline
\textbf{Code} & \textbf{Meaning} \\
\hline
\textbf{200} & OK - Request succeeded \\
\hline
\textbf{201} & Created - Resource created successfully \\
\hline
\textbf{400} & Bad Request - Invalid input, missing required fields \\
\hline
\textbf{401} & Unauthorized - Missing or invalid token \\
\hline
\textbf{403} & Forbidden - Insufficient permissions for this resource \\
\hline
\textbf{404} & Not Found - Resource doesn't exist \\
\hline
\textbf{409} & Conflict - Resource already exists (e.g., duplicate email) \\
\hline
\textbf{422} & Unprocessable Entity - Validation error (see detail field) \\
\hline
\textbf{500} & Server Error - Internal server error \\
\hline
\textbf{503} & Service Unavailable - Server temporarily down \\
\hline
\end{tabular}
\caption{HTTP Status Codes Reference}
\end{table}

\newpage

% ============================================================================
% SECTION 4: ERROR RESPONSES
% ============================================================================
\section{Error Handling}

All error responses follow a consistent format:

\begin{lstlisting}
{
  "detail": "Human-readable error message",
  "error_code": "ERROR_CODE_ENUM",
  "timestamp": "2026-05-01T11:30:00Z"
}
\end{lstlisting}

\subsection{Common Errors}

\begin{tcolorbox}[colback=lightred, colframe=darkblue]
\textbf{401 Unauthorized}
\begin{lstlisting}
{
  "detail": "Not authenticated",
  "error_code": "UNAUTHENTICATED"
}
\end{lstlisting}
\textbf{Action}: Include valid Bearer token in Authorization header.
\end{tcolorbox}

\begin{tcolorbox}[colback=lightred, colframe=darkblue]
\textbf{403 Forbidden}
\begin{lstlisting}
{
  "detail": "Insufficient permissions",
  "error_code": "INSUFFICIENT_PERMISSIONS"
}
\end{lstlisting}
\textbf{Action}: Verify user role or contact admin.
\end{tcolorbox}

\begin{tcolorbox}[colback=lightred, colframe=darkblue]
\textbf{400 Bad Request}
\begin{lstlisting}
{
  "detail": "Invalid input provided",
  "error_code": "VALIDATION_ERROR"
}
\end{lstlisting}
\textbf{Action}: Check request format and required fields.
\end{tcolorbox}

\newpage

% ============================================================================
% SECTION 5: AUTHENTICATION ENDPOINTS
% ============================================================================
\section{Authentication Endpoints}

\subsection{POST /auth/login}

Login user and receive authentication tokens.

\subsubsection{Request}

\begin{lstlisting}
POST /api/v1/auth/login
Content-Type: application/json

{
  "email": "user@visiolearn.org",
  "password": "SecurePass123!"
}
\end{lstlisting}

\subsubsection{Response (200 OK)}

\begin{lstlisting}
{
  "access_token": "eyJhbGciOiJIUzI1NiIsInR5cCI6IkpXVCJ9...",
  "refresh_token": "eyJhbGciOiJIUzI1NiIsInR5cCI6IkpXVCJ9...",
  "token_type": "bearer",
  "expires_in": 3600,
  "user": {
    "id": "550e8400-e29b-41d4-a716-446655440000",
    "email": "user@visiolearn.org",
    "role": "student"
  }
}
\end{lstlisting}

\subsubsection{Error Responses}

\begin{table}[H]
\centering
\small
\begin{tabular}{|l|l|}
\hline
\textbf{Status} & \textbf{Reason} \\
\hline
401 & Invalid email or password \\
\hline
400 & Missing email or password field \\
\hline
\end{tabular}
\end{table}

\subsection{POST /auth/refresh}

Get a new access token using refresh token.

\subsubsection{Request}

\begin{lstlisting}
POST /api/v1/auth/refresh
Authorization: Bearer <refresh_token>
\end{lstlisting}

\subsubsection{Response (200 OK)}

\begin{lstlisting}
{
  "access_token": "eyJhbGciOiJIUzI1NiIsInR5cCI6IkpXVCJ9...",
  "refresh_token": "eyJhbGciOiJIUzI1NiIsInR5cCI6IkpXVCJ9...",
  "token_type": "bearer",
  "expires_in": 3600
}
\end{lstlisting}

\textbf{Note}: The old refresh token is invalidated and a new one is returned. Store the new refresh token immediately.

\subsection{POST /auth/logout}

Logout user and invalidate tokens.

\subsubsection{Request}

\begin{lstlisting}
POST /api/v1/auth/logout
Authorization: Bearer <access_token>
\end{lstlisting}

\subsubsection{Response (200 OK)}

\begin{lstlisting}
{
  "message": "Successfully logged out"
}
\end{lstlisting}

\newpage

% ============================================================================
% SECTION 6: USER ENDPOINTS
% ============================================================================
\section{User Management Endpoints}

All user endpoints require authentication.

\subsection{GET /users/me}

Get current logged-in user's profile.

\subsubsection{Request}

\begin{lstlisting}
GET /api/v1/users/me
Authorization: Bearer <access_token>
\end{lstlisting}

\subsubsection{Response (200 OK)}

\begin{lstlisting}
{
  "id": "550e8400-e29b-41d4-a716-446655440000",
  "email": "student@visiolearn.org",
  "role": "student",
  "school_id": "f5c8f3b1-2a4c-4d6e-8f0a-1b3c5d7e9f1a",
  "created_at": "2026-05-01T10:00:00Z"
}
\end{lstlisting}

\subsection{GET /users}

Get list of all users (admin only).

\subsubsection{Request}

\begin{lstlisting}
GET /api/v1/users
Authorization: Bearer <admin_access_token>
\end{lstlisting}

\subsubsection{Query Parameters}

\begin{table}[H]
\centering
\small
\begin{tabular}{|l|l|l|}
\hline
\textbf{Parameter} & \textbf{Type} & \textbf{Description} \\
\hline
\texttt{skip} & integer & Pagination offset (default: 0) \\
\hline
\texttt{limit} & integer & Records per page (default: 10, max: 100) \\
\hline
\texttt{role} & string & Filter by role (admin, teacher, student) \\
\hline
\end{tabular}
\end{table}

\subsubsection{Response (200 OK)}

\begin{lstlisting}
{
  "users": [
    {
      "id": "550e8400-e29b-41d4-a716-446655440000",
      "email": "student@visiolearn.org",
      "role": "student",
      "school_id": "f5c8f3b1-2a4c-4d6e-8f0a-1b3c5d7e9f1a",
      "created_at": "2026-05-01T10:00:00Z"
    }
  ],
  "total": 150,
  "skip": 0,
  "limit": 10
}
\end{lstlisting}

\newpage

% ============================================================================
% SECTION 7: NOTES ENDPOINTS
% ============================================================================
\section{Lesson Notes Endpoints}

Upload and manage lesson materials (PDF, DOCX, TXT).

\subsection{POST /notes}

Upload a new lesson note.

\subsubsection{Requirements}

\begin{itemize}
    \item \textbf{Authentication}: Required
    \item \textbf{Role}: teacher or admin only
    \item \textbf{File Size}: Max 50 MB
    \item \textbf{File Types}: PDF, DOCX, TXT only
    \item \textbf{Content-Type}: multipart/form-data
\end{itemize}

\subsubsection{Request}

\begin{lstlisting}
POST /api/v1/notes
Authorization: Bearer <teacher_access_token>
Content-Type: multipart/form-data

form-data:
  - title: "Chapter 1: Introduction"
  - description: "Introduction to Python basics"
  - file: <binary_file>
  - unit_id: "308685dc-7b35-4494-b62b-f9ee2de4b884"
  - subject: "Python"
\end{lstlisting}

\subsubsection{Form Parameters}

\begin{table}[H]
\centering
\small
\begin{tabular}{|l|l|l|l|}
\hline
\textbf{Field} & \textbf{Type} & \textbf{Required} & \textbf{Description} \\
\hline
\texttt{title} & string & Yes & Note title (max 255 chars) \\
\hline
\texttt{description} & string & No & Detailed description \\
\hline
\texttt{file} & binary & Yes & PDF/DOCX/TXT file \\
\hline
\texttt{unit\_id} & UUID & Yes & Associated unit ID \\
\hline
\texttt{subject} & string & No & Subject area \\
\hline
\end{tabular}
\end{table}

\subsubsection{Response (201 Created)}

\begin{lstlisting}
{
  "id": "a1b2c3d4-e5f6-47g8-h9i0-j1k2l3m4n5o6",
  "title": "Chapter 1: Introduction",
  "description": "Introduction to Python basics",
  "unit_id": "308685dc-7b35-4494-b62b-f9ee2de4b884",
  "file_url": "https://storage.visiolearn.com/notes/a1b2c3d4.pdf",
  "uploaded_by": "550e8400-e29b-41d4-a716-446655440000",
  "created_at": "2026-05-01T11:30:00Z",
  "status": "processing"
}
\end{lstlisting}

\subsubsection{Error Responses}

\begin{tcolorbox}[colback=lightred, colframe=darkblue]
\textbf{400 - File too large or wrong format}
\begin{lstlisting}
{
  "detail": "File must be PDF, DOCX, or TXT and under 50 MB"
}
\end{lstlisting}
\end{tcolorbox}

\begin{tcolorbox}[colback=lightred, colframe=darkblue]
\textbf{403 - Insufficient permissions}
\begin{lstlisting}
{
  "detail": "Only teachers can upload notes"
}
\end{lstlisting}
\end{tcolorbox}

\subsection{GET /notes}

Get list of lesson notes (admin/teacher) or assigned notes (student).

\subsubsection{Request}

\begin{lstlisting}
GET /api/v1/notes?skip=0&limit=10&status=completed
Authorization: Bearer <access_token>
\end{lstlisting}

\subsubsection{Query Parameters}

\begin{table}[H]
\centering
\small
\begin{tabular}{|l|l|l|}
\hline
\textbf{Parameter} & \textbf{Type} & \textbf{Description} \\
\hline
\texttt{skip} & integer & Pagination offset \\
\hline
\texttt{limit} & integer & Records per page (max: 100) \\
\hline
\texttt{status} & string & Filter: processing, completed, failed \\
\hline
\texttt{unit\_id} & UUID & Filter by unit \\
\hline
\end{tabular}
\end{table}

\subsubsection{Response (200 OK)}

\begin{lstlisting}
{
  "notes": [
    {
      "id": "a1b2c3d4-e5f6-47g8-h9i0-j1k2l3m4n5o6",
      "title": "Chapter 1: Introduction",
      "description": "Introduction to Python basics",
      "unit_id": "308685dc-7b35-4494-b62b-f9ee2de4b884",
      "file_url": "https://storage.visiolearn.com/notes/a1b2c3d4.pdf",
      "status": "completed",
      "created_at": "2026-05-01T11:30:00Z"
    }
  ],
  "total": 45,
  "skip": 0,
  "limit": 10
}
\end{lstlisting}

\subsection{GET /notes/\{note\_id\}}

Get details of a specific note.

\subsubsection{Request}

\begin{lstlisting}
GET /api/v1/notes/a1b2c3d4-e5f6-47g8-h9i0-j1k2l3m4n5o6
Authorization: Bearer <access_token>
\end{lstlisting}

\subsubsection{Response (200 OK)}

\begin{lstlisting}
{
  "id": "a1b2c3d4-e5f6-47g8-h9i0-j1k2l3m4n5o6",
  "title": "Chapter 1: Introduction",
  "description": "Introduction to Python basics",
  "unit_id": "308685dc-7b35-4494-b62b-f9ee2de4b884",
  "file_url": "https://storage.visiolearn.com/notes/a1b2c3d4.pdf",
  "uploaded_by": "550e8400-e29b-41d4-a716-446655440000",
  "status": "completed",
  "created_at": "2026-05-01T11:30:00Z",
  "updated_at": "2026-05-01T12:00:00Z"
}
\end{lstlisting}

\newpage

% ============================================================================
% SECTION 8: VOICE SESSION ENDPOINTS
% ============================================================================
\section{Voice Session Endpoints}

Manage voice-based learning sessions where students interact with AI.

\subsection{Understanding Voice Session Lifecycle}

Voice sessions follow a strict lifecycle:

\begin{enumerate}
    \item \textbf{ACTIVE}: Session started, student can send voice/text input
    \item \textbf{PAUSED}: Student paused the session temporarily
    \item \textbf{COMPLETED}: Session ended, results are finalized
\end{enumerate}

Once a session is \textbf{COMPLETED}, it cannot be modified.

\subsection{POST /voice/sessions}

Start a new voice learning session.

\subsubsection{Requirements}

\begin{itemize}
    \item \textbf{Authentication}: Required
    \item \textbf{Role}: student only
    \item \textbf{note\_id}: Must be a valid note that the student has access to
\end{itemize}

\subsubsection{Request}

\begin{lstlisting}
POST /api/v1/voice/sessions
Authorization: Bearer <student_access_token>
Content-Type: application/json

{
  "student_id": "550e8400-e29b-41d4-a716-446655440000",
  "note_id": "a1b2c3d4-e5f6-47g8-h9i0-j1k2l3m4n5o6",
  "unit_id": "308685dc-7b35-4494-b62b-f9ee2de4b884"
}
\end{lstlisting}

\subsubsection{Request Fields}

\begin{table}[H]
\centering
\small
\begin{tabular}{|l|l|l|}
\hline
\textbf{Field} & \textbf{Type} & \textbf{Description} \\
\hline
\texttt{student\_id} & UUID & Your user ID \\
\hline
\texttt{note\_id} & UUID & ID of lesson note to study \\
\hline
\texttt{unit\_id} & UUID & ID of course unit \\
\hline
\end{tabular}
\end{table}

\subsubsection{Response (201 Created)}

\begin{lstlisting}
{
  "id": "sess_1a2b3c4d5e6f7g8h",
  "student_id": "550e8400-e29b-41d4-a716-446655440000",
  "note_id": "a1b2c3d4-e5f6-47g8-h9i0-j1k2l3m4n5o6",
  "unit_id": "308685dc-7b35-4494-b62b-f9ee2de4b884",
  "status": "ACTIVE",
  "started_at": "2026-05-01T11:30:00Z",
  "duration_seconds": 0,
  "interactions_count": 0
}
\end{lstlisting}

\subsubsection{Error Responses}

\begin{tcolorbox}[colback=lightred, colframe=darkblue]
\textbf{403 - Unauthorized Student}
\begin{lstlisting}
{
  "detail": "Unauthorized: student_id must match current user"
}
\end{lstlisting}
\end{tcolorbox}

\subsection{POST /voice/sessions/\{session\_id\}/interactions}

Add interaction (question/answer) to an active session.

\subsubsection{Request}

\begin{lstlisting}
POST /api/v1/voice/sessions/sess_1a2b3c4d5e6f7g8h/interactions
Authorization: Bearer <student_access_token>
Content-Type: application/json

{
  "input_type": "voice",
  "input_text": "What is Python?",
  "ai_response": "Python is a programming language..."
}
\end{lstlisting}

\subsubsection{Request Fields}

\begin{table}[H]
\centering
\small
\begin{tabular}{|l|l|l|}
\hline
\textbf{Field} & \textbf{Type} & \textbf{Description} \\
\hline
\texttt{input\_type} & string & voice or text \\
\hline
\texttt{input\_text} & string & Student's question \\
\hline
\texttt{ai\_response} & string & AI's response \\
\hline
\end{tabular}
\end{table}

\subsubsection{Response (201 Created)}

\begin{lstlisting}
{
  "id": "int_1a2b3c4d5e6f7g8h",
  "session_id": "sess_1a2b3c4d5e6f7g8h",
  "sequence_number": 1,
  "input_type": "voice",
  "input_text": "What is Python?",
  "ai_response": "Python is a programming language...",
  "created_at": "2026-05-01T11:30:15Z"
}
\end{lstlisting}

\subsection{PUT /voice/sessions/\{session\_id\}}

Update session status (pause, resume, complete).

\subsubsection{Request}

\begin{lstlisting}
PUT /api/v1/voice/sessions/sess_1a2b3c4d5e6f7g8h
Authorization: Bearer <student_access_token>
Content-Type: application/json

{
  "status": "COMPLETED"
}
\end{lstlisting}

\subsubsection{Valid Status Transitions}

\begin{table}[H]
\centering
\small
\begin{tabular}{|l|l|}
\hline
\textbf{Current Status} & \textbf{Valid Next Status} \\
\hline
ACTIVE & PAUSED, COMPLETED \\
\hline
PAUSED & ACTIVE, COMPLETED \\
\hline
COMPLETED & (no transitions allowed) \\
\hline
\end{tabular}
\end{table}

\subsubsection{Response (200 OK)}

\begin{lstlisting}
{
  "id": "sess_1a2b3c4d5e6f7g8h",
  "student_id": "550e8400-e29b-41d4-a716-446655440000",
  "note_id": "a1b2c3d4-e5f6-47g8-h9i0-j1k2l3m4n5o6",
  "unit_id": "308685dc-7b35-4494-b62b-f9ee2de4b884",
  "status": "COMPLETED",
  "started_at": "2026-05-01T11:30:00Z",
  "ended_at": "2026-05-01T11:45:30Z",
  "duration_seconds": 930,
  "interactions_count": 5
}
\end{lstlisting}

\subsection{GET /voice/sessions}

Get list of user's voice sessions.

\subsubsection{Request}

\begin{lstlisting}
GET /api/v1/voice/sessions?status=COMPLETED&skip=0&limit=10
Authorization: Bearer <access_token>
\end{lstlisting}

\subsubsection{Query Parameters}

\begin{table}[H]
\centering
\small
\begin{tabular}{|l|l|l|}
\hline
\textbf{Parameter} & \textbf{Type} & \textbf{Description} \\
\hline
\texttt{status} & string & Filter: ACTIVE, PAUSED, COMPLETED \\
\hline
\texttt{skip} & integer & Pagination offset \\
\hline
\texttt{limit} & integer & Records per page (max: 100) \\
\hline
\end{tabular}
\end{table}

\subsubsection{Response (200 OK)}

\begin{lstlisting}
{
  "sessions": [
    {
      "id": "sess_1a2b3c4d5e6f7g8h",
      "student_id": "550e8400-e29b-41d4-a716-446655440000",
      "note_id": "a1b2c3d4-e5f6-47g8-h9i0-j1k2l3m4n5o6",
      "status": "COMPLETED",
      "started_at": "2026-05-01T11:30:00Z",
      "ended_at": "2026-05-01T11:45:30Z",
      "duration_seconds": 930,
      "interactions_count": 5
    }
  ],
  "total": 12,
  "skip": 0,
  "limit": 10
}
\end{lstlisting}

\subsection{GET /voice/sessions/\{session\_id\}}

Get details of a specific voice session.

\subsubsection{Request}

\begin{lstlisting}
GET /api/v1/voice/sessions/sess_1a2b3c4d5e6f7g8h
Authorization: Bearer <access_token>
\end{lstlisting}

\subsubsection{Response (200 OK)}

\begin{lstlisting}
{
  "id": "sess_1a2b3c4d5e6f7g8h",
  "student_id": "550e8400-e29b-41d4-a716-446655440000",
  "note_id": "a1b2c3d4-e5f6-47g8-h9i0-j1k2l3m4n5o6",
  "unit_id": "308685dc-7b35-4494-b62b-f9ee2de4b884",
  "status": "COMPLETED",
  "started_at": "2026-05-01T11:30:00Z",
  "ended_at": "2026-05-01T11:45:30Z",
  "duration_seconds": 930,
  "interactions_count": 5,
  "interactions": [
    {
      "id": "int_1a2b3c4d5e6f7g8h",
      "sequence_number": 1,
      "input_type": "voice",
      "input_text": "What is Python?",
      "ai_response": "Python is a programming language...",
      "created_at": "2026-05-01T11:30:15Z"
    }
  ]
}
\end{lstlisting}

\newpage

% ============================================================================
% SECTION 9: COMMON PATTERNS & BEST PRACTICES
% ============================================================================
\section{Common Integration Patterns}

\subsection{Token Refresh Flow}

When you receive a 401 response:

\begin{enumerate}
    \item Use your \texttt{refresh\_token} to call \texttt{POST /auth/refresh}
    \item Store the new \texttt{access\_token}
    \item Retry the original failed request with the new token
\end{enumerate}

\textbf{Important}: The refresh token itself is also rotated, so always store the new one returned.

\subsection{Pagination}

All list endpoints support pagination:

\begin{lstlisting}
GET /api/v1/notes?skip=20&limit=10
\end{lstlisting}

\begin{table}[H]
\centering
\small
\begin{tabular}{|l|l|}
\hline
\textbf{Parameter} & \textbf{Meaning} \\
\hline
\texttt{skip} & Number of records to skip (offset) \\
\hline
\texttt{limit} & Number of records to return \\
\hline
\texttt{total} & Total available records \\
\hline
\end{tabular}
\end{table}

\subsection{Handling File Uploads}

For file uploads (lesson notes):

\begin{enumerate}
    \item Use \texttt{Content-Type: multipart/form-data}
    \item \textbf{Do NOT} set Content-Type header manually—let your HTTP client set it automatically with boundary
    \item Include all required form fields (title, file, unit\_id)
    \item Max file size: 50 MB
    \item Supported formats: PDF, DOCX, TXT
\end{enumerate}

\subsection{Voice Session Flow}

Typical student workflow:

\begin{enumerate}
    \item \textbf{Start}: \texttt{POST /voice/sessions} with note\_id
    \item \textbf{Interact}: \texttt{POST /voice/sessions/\{id\}/interactions} repeatedly
    \item \textbf{Pause}: \texttt{PUT /voice/sessions/\{id\}} with \texttt{status: PAUSED}
    \item \textbf{Resume}: \texttt{PUT /voice/sessions/\{id\}} with \texttt{status: ACTIVE}
    \item \textbf{End}: \texttt{PUT /voice/sessions/\{id\}} with \texttt{status: COMPLETED}
\end{enumerate}

Once COMPLETED, the session becomes read-only.

\newpage

% ============================================================================
% SECTION 10: QUICK REFERENCE
% ============================================================================
\section{Quick API Reference}

\begin{table}[H]
\centering
\small
\begin{longtable}{|l|l|l|l|}
\hline
\textbf{Method} & \textbf{Endpoint} & \textbf{Auth} & \textbf{Purpose} \\
\hline
\endhead

POST & /auth/login & No & Login user \\
\hline
POST & /auth/refresh & Refresh Token & Get new access token \\
\hline
POST & /auth/logout & Yes & Logout user \\
\hline
GET & /users/me & Yes & Get current user \\
\hline
GET & /users & Admin & List all users \\
\hline
POST & /notes & Teacher & Upload lesson note \\
\hline
GET & /notes & Yes & List notes \\
\hline
GET & /notes/\{id\} & Yes & Get note details \\
\hline
POST & /voice/sessions & Student & Start voice session \\
\hline
GET & /voice/sessions & Yes & List voice sessions \\
\hline
GET & /voice/sessions/\{id\} & Yes & Get session details \\
\hline
POST & /voice/sessions/\{id\}/interactions & Student & Add interaction \\
\hline
PUT & /voice/sessions/\{id\} & Student & Update session status \\
\hline

\caption{Complete Endpoint Reference}
\end{longtable}
\end{table}

\newpage

% ============================================================================
% SECTION 11: DATA FORMATS & VALIDATION
% ============================================================================
\section{Data Formats \& Validation}

\subsection{Email Format}

All email fields must be valid RFC 5322 format:

\begin{lstlisting}
Valid: user@visiolearn.org, teacher.name@example.com
Invalid: user@, user@.com, user@example
\end{lstlisting}

\subsection{Password Requirements}

Passwords must meet these requirements:

\begin{itemize}
    \item Minimum 8 characters
    \item At least one uppercase letter (A-Z)
    \item At least one lowercase letter (a-z)
    \item At least one number (0-9)
    \item At least one special character (!@\#\$\%\^\&*)
\end{itemize}

Example: \texttt{SecurePass123!}

\subsection{UUID Format}

All IDs use UUID v4 format:

\begin{lstlisting}
550e8400-e29b-41d4-a716-446655440000
\end{lstlisting}

\subsection{Timestamp Format}

All timestamps are in ISO 8601 format (UTC):

\begin{lstlisting}
2026-05-01T11:30:00Z
\end{lstlisting}

\subsection{File Upload Constraints}

\begin{table}[H]
\centering
\small
\begin{tabular}{|l|l|}
\hline
\textbf{Constraint} & \textbf{Value} \\
\hline
Max file size & 50 MB \\
\hline
Allowed formats & PDF, DOCX, TXT \\
\hline
Encoding type & multipart/form-data \\
\hline
\end{tabular}
\end{table}

\newpage

% ============================================================================
% SECTION 12: TROUBLESHOOTING
% ============================================================================
\section{Troubleshooting Guide}

\subsection{401 Unauthorized}

\begin{tcolorbox}[colback=lightred, colframe=darkblue]
\textbf{Problem}: Request returns 401 Unauthorized

\textbf{Causes}:
\begin{itemize}
    \item Missing Authorization header
    \item Invalid or expired access token
    \item Malformed Bearer token format
\end{itemize}

\textbf{Solution}:
\begin{itemize}
    \item Verify Authorization header is present: \texttt{Authorization: Bearer <token>}
    \item Check token hasn't expired (expires\_in: 3600 seconds)
    \item Use \texttt{POST /auth/refresh} to get a new token
    \item Re-login if refresh fails
\end{itemize}
\end{tcolorbox}

\subsection{403 Forbidden}

\begin{tcolorbox}[colback=lightred, colframe=darkblue]
\textbf{Problem}: Request returns 403 Forbidden

\textbf{Causes}:
\begin{itemize}
    \item User role doesn't have permission for endpoint
    \item Student trying to upload notes (teacher only)
    \item Teacher accessing admin-only endpoints
\end{itemize}

\textbf{Solution}:
\begin{itemize}
    \item Check endpoint role requirements
    \item Use correct user account (teacher can't upload as student)
    \item Contact admin if you need elevated permissions
\end{itemize}
\end{tcolorbox}

\subsection{400 Bad Request}

\begin{tcolorbox}[colback=lightred, colframe=darkblue]
\textbf{Problem}: Request returns 400 Bad Request

\textbf{Causes}:
\begin{itemize}
    \item Missing required field
    \item Invalid data format
    \item Invalid UUID or email format
\end{itemize}

\textbf{Solution}:
\begin{itemize}
    \item Check all required fields are provided
    \item Validate email and UUID formats before sending
    \item Verify JSON syntax is correct
    \item Check Content-Type header matches request body
\end{itemize}
\end{tcolorbox}

\subsection{422 Unprocessable Entity}

\begin{tcolorbox}[colback=lightred, colframe=darkblue]
\textbf{Problem}: Request returns 422 Unprocessable Entity

\textbf{Causes}:
\begin{itemize}
    \item Validation error on field values
    \item Field constraints not met (e.g., password too weak)
    \item Type mismatch (string vs number)
\end{itemize}

\textbf{Solution}:
\begin{itemize}
    \item Check the \texttt{detail} field for specific validation error
    \item Verify password meets complexity requirements
    \item Ensure data types match API schema
    \item Re-read data format section above
\end{itemize}
\end{tcolorbox}

\subsection{500 Server Error}

\begin{tcolorbox}[colback=lightred, colframe=darkblue]
\textbf{Problem}: Request returns 500 Server Error

\textbf{Causes}:
\begin{itemize}
    \item Server-side issue or bug
    \item Database connection problem
    \item Service temporarily unavailable
\end{itemize}

\textbf{Solution}:
\begin{itemize}
    \item Retry the request after a few seconds
    \item Check status page: https://visiolearn-backend.onrender.com/health
    \item Contact backend team if error persists
\end{itemize}
\end{tcolorbox}

\newpage

% ============================================================================
% SECTION 13: USEFUL LINKS
% ============================================================================
\section{Resources}

\subsection{Live API Documentation}

\begin{itemize}
    \item \textbf{Swagger UI}: \url{https://visiolearn-backend.onrender.com/docs}
    \item \textbf{ReDoc}: \url{https://visiolearn-backend.onrender.com/redoc}
    \item \textbf{OpenAPI JSON}: \url{https://visiolearn-backend.onrender.com/openapi.json}
\end{itemize}

\subsection{Testing}

Use these tools to test API endpoints:

\begin{itemize}
    \item \textbf{Postman}: \url{https://www.postman.com} (GUI-based)
    \item \textbf{cURL}: Command-line tool (pre-installed on most systems)
    \item \textbf{Insomnia}: \url{https://insomnia.rest} (open-source alternative)
\end{itemize}

\subsection{HTTP Clients in Popular Languages}

\begin{table}[H]
\centering
\small
\begin{tabular}{|l|l|}
\hline
\textbf{Language} & \textbf{HTTP Client} \\
\hline
JavaScript & Axios, Fetch API, node-fetch \\
\hline
Python & requests, httpx \\
\hline
Java & OkHttp, HttpClient \\
\hline
Go & net/http, resty \\
\hline
C\# & HttpClient, RestSharp \\
\hline
\end{tabular}
\end{table}

\subsection{Support}

\begin{itemize}
    \item \textbf{GitHub Issues}: \url{https://github.com/Nyanja23/VisioLearn-Backend/issues}
    \item \textbf{Swagger UI}: Try endpoints live at \url{https://visiolearn-backend.onrender.com/docs}
\end{itemize}

\end{document}
